\section{Research proposals}\label{sec:future}

\subsection{Preliminaries: a working definition of failure attribution}\label{sec:future:prelim}

\begin{multicols}{2}
\paragraph{Why a definition is needed} The three proposals that follow need a target that stays fixed from one dataset to the next. The survey shows that the current targets do not:
\begin{itemize}[leftmargin=1.1em]
\item Every dataset defines its labelled step \emph{counterfactually}---``the earliest step such that, had it been correct, the run would have succeeded''---but only TraceElephant ever tests the counterfactual by re-execution (\S\ref{ds:traceelephant}).
\item The summary that is labelled differs: the earliest decisive step (Who\&When, TrajErrBench, LongRCA), the first \emph{unrecovered} failure (AgentRx), a set of steps with per-annotator votes (MP-Bench), every error span (TRAIL), three timestamps $t_{\mathrm{err}}/t_{\mathrm{lock}}/t_{\mathrm{obs}}$ (Failure-as-a-Process), a step with a tolerance span (SearchAuditBench) or a set of (edge, fault side, mode) triples (MegaRCA-Mix).
\item Humans do not agree on the proxy: exact-step agreement is 40.3\% in LongRCA, 77.9\% in AgentHallu and 91.2\% / 66.7\% on the two TrajErrBench domains, TraceElephant's round-1 Krippendorff's $\alpha$ for the step is 0.64, and the blamed agent agrees between Who\&When and TraceElephant on only 18 of 86 paired tasks.
\end{itemize}
We therefore take the counterfactual itself as the definition and treat every human label as a noisy estimate of it.

\paragraph{Setting and notation} An agent system $\pi$ (single- or multi-agent, stochastic) interacts with an environment $\mathcal{E}$ on a task $x$ with an automatic success predicate $S\in\{0,1\}$ (ground-truth answer, unit tests, verifier). A trajectory is $\tau=(s_1,\dots,s_T)$ where step $s_t=(a_t,u_t,o_t)$ records the acting agent $a_t\in\mathcal{A}$, its output $u_t$ (message, tool call or code) and the resulting observation $o_t$ (tool result, page, test log). $\tau_{<k}$ denotes the prefix $(s_1,\dots,s_{k-1})$. We write $\mathrm{Roll}(\tau_{<k}\oplus s_k')$ for the random continuation obtained by restoring the state reached after $\tau_{<k}$, executing the (possibly modified) step $s_k'$ and letting $\pi$ and $\mathcal{E}$ run to termination; randomness comes from sampling in $\pi$ and from $\mathcal{E}$.

\begin{assumption}[Replayability]\label{as:replay}
The environment state after $\tau_{<k}$ can be restored exactly (container snapshot, sandbox seed, simulated user), so that $\mathrm{Roll}(\cdot)$ is a well-defined distribution over continuations. Live-web environments violate this: SymTrace's selective replay \cite{luan2026symtrace} serves the prefix from recorded tool returns, which restores the prefix exactly but leaves the continuation unreproducible (and the authors note that temperature-0 calls to a hosted model do not guarantee determinism either), so live-web cases are excluded from the measured part of the programme.
\end{assumption}

\begin{assumption}[Fixed intervention vocabulary]\label{as:vocab}
Interventions come from a finite, declared set $\Delta=\Delta_{\mathrm{fix}}\cup\Delta_{\mathrm{inj}}$; each $\delta\in\Delta$ is a deterministic map $s_k\mapsto\delta(s_k)$ that changes exactly one step. Table~\ref{tab:vocab} lists a starting vocabulary. ``A change to a step'' without such a vocabulary is not a learnable object.
\end{assumption}

\begin{definition}[Intervention effect]\label{def:effect}
For a trajectory $\tau$, a step index $k$ and an intervention $\delta\in\Delta$, define the \emph{baseline success rate} and the \emph{intervened success rate}
\begin{align*}
p_0(\tau,k)&=\Pr\big[S(\mathrm{Roll}(\tau_{<k}\oplus s_k))=1\big],\\
p_\delta(\tau,k)&=\Pr\big[S(\mathrm{Roll}(\tau_{<k}\oplus\delta(s_k)))=1\big],
\end{align*}
and the \emph{effect} $\phi(\tau,k,\delta)=p_\delta(\tau,k)-p_0(\tau,k)\in[-1,1]$. Both probabilities are estimated by $N$ independent rollouts; $p_0$ is estimated with the original step re-executed (not read from the observed outcome), which is what makes the quantity robust to run-to-run stochasticity. A step $k$ is \emph{$\alpha$-decisive} for a failed $\tau$ if $\max_{\delta\in\Delta_{\mathrm{fix}}}\phi(\tau,k,\delta)\ge\alpha$, and \emph{$\alpha$-fragile} for a successful $\tau$ if $\min_{\delta\in\Delta_{\mathrm{inj}}}\phi(\tau,k,\delta)\le-\alpha$.
\end{definition}

\begin{definition}[Failure attribution]\label{def:fa}
For a failed trajectory $\tau$ the \emph{repair-effect profile} is
\[e_\tau(k)=\max_{\delta\in\Delta_{\mathrm{fix}}}\phi(\tau,k,\delta),\qquad k=1,\dots,T,\]
and \emph{failure attribution} is the estimation of $e_\tau$. For a threshold $\alpha\in(0,1]$ that is stated with every result, the labels used in the literature are summary statistics of the profile:
\begin{itemize}[leftmargin=1.1em]
\item \textbf{error set} $E_\alpha(\tau)=\{k:\ e_\tau(k)\ge\alpha\}$ --- the steps whose repair would have rescued the run;
\item \textbf{decisive step} $k^\star=\min E_\alpha(\tau)$ --- the earliest of them;
\item \textbf{lock-in} $k_{\mathrm{lock}}=\max E_\alpha(\tau)+1$ --- the first step from which no single-step repair recovers the run;
\item \textbf{responsible agent} $a^\star=a_{k^\star}$ --- the agent that emitted the decisive step; no separate ``who'' label is defined;
\item \textbf{fault side} --- whether the maximising $\delta$ at $k^\star$ acts on the output $u_k$ (model-side) or on the observation $o_k$ (tool- or environment-side).
\end{itemize}
\end{definition}

\paragraph{Correspondence with the surveyed labels}
\begin{itemize}[leftmargin=1.1em]
\item $k^\star$ is the target of Who\&When, TraceElephant, TrajErrBench, LongRCA, CORRECT-Error and the injected datasets, and the $t_{\mathrm{err}}$ of Failure-as-a-Process.
\item $E_\alpha$ is what MP-Bench's vote counts, AgentRx's failure events and TRAIL's error spans approximate; MP-Bench's per-step consensus rate is a human estimate of $e_\tau(k)$ itself.
\item $k_{\mathrm{lock}}$ is AgentRx's ``first failure from which the run never recovers'' and Failure-as-a-Process's $t_{\mathrm{lock}}$; $k_{\mathrm{lock}}-k^\star$ is that paper's fix window (median 1 step).
\item $a^\star$: the sender of the labelled step equals the responsible role in 99.2\% of LongRCA and 100\% of TracerTraj records, so defining the agent through the step loses nothing and removes the naming problems of Section~\ref{sec:reco}.
\item Fault side recovers MegaRCA-Mix's model / harness / environment / grader sides and the environment-caused failures that SearchAuditBench excludes.
\end{itemize}

\paragraph{What the definition is robust to}
\begin{itemize}[leftmargin=1.1em]
\item \emph{Annotator variance.} The target is measured, not adjudicated; a dataset whose ground truth depends on who labelled it cannot be a stable target across datasets.
\item \emph{Run-to-run stochasticity.} $p_0$ is re-estimated by re-execution, so a run that failed by chance has a small effect everywhere instead of a forced label.
\item \emph{Several causes.} The error set is graded by effect size; MP-Bench's 16\% single-step agreement becomes a property of the profile (several steps of comparable effect) rather than label noise.
\item \emph{Recovery.} Decisive step and lock-in come from one profile, so the two conventions in the literature are reconciled instead of chosen between.
\item \emph{Observation level.} The definition lives on the system, not on the log: output-only logs, LLM-call traces and span trees are different estimators of the same quantity, not different tasks.
\end{itemize}

\paragraph{What must be fixed for it to be well posed}
\begin{itemize}[leftmargin=1.1em]
\item \emph{The vocabulary $\Delta$} (Assumption~\ref{as:vocab}, Table~\ref{tab:vocab}). ``Had it been correct'' is not learnable until ``correct'' is an operation; results are reported per intervention. Errors with no single-step intervention---a wrong team composition before the log begins, which is where 20 of 126 Who\&When Algorithm-Generated labels sit (step~0)---are declared out of scope, not scored as zero.
\item \emph{The step unit.} One step is one policy decision: an agent output $u_t$ together with the observation $o_t$ it produced. Orchestrator bookkeeping, user turns and bare tool returns are context, not candidates. This removes the indexing differences catalogued in Section~\ref{sec:reco}.
\item \emph{$N$ and $\alpha$.} Effects are reported with intervals; $N=5$--$10$ rollouts separate an effect of 0.5 from 0 but not 0.1 from 0, so $\alpha$ is not set below the resolution that $N$ affords.
\end{itemize}

\paragraph{Proxies when replay is impossible}
\begin{itemize}[leftmargin=1.1em]
\item Live-web GAIA and AssistantBench runs and output-only logs violate Assumption~\ref{as:replay}. There the human decisive step is used as a proxy for $k^\star$, scored with exact match \emph{and} a $\pm k$ window, and the agreement figures above are quoted as its ceiling.
\item No method is evaluated on proxies alone: the replayable environments of the survey (SWE-bench and SWE-Bench~Pro, ALFWorld, WebShop, $\tau^2$-Bench, TravelPlanner, VitaBench, Terminal-Bench) carry the measured version; the proxies test transfer.
\item In one sentence: failure attribution is the estimation of a step-level repair-effect profile under a declared intervention set, and the decisive step, the error set, the lock-in point and the responsible agent are summary statistics of that profile.
\end{itemize}
\end{multicols}

\medskip\noindent\begin{minipage}{\linewidth}\centering\small
\begin{tabularx}{\linewidth}{@{}l l >{\RaggedRight}X@{}}\toprule
Set & Intervention $\delta$ & Definition (changes step $k$ only) \\\midrule
$\Delta_{\mathrm{fix}}$ & \textsc{oracle-action} & replace $u_k$ by the action a human/LLM with access to the ground truth deems correct at $k$ \\
 & \textsc{hint} & prepend to $u_k$'s context a one-sentence statement of the fact the step got wrong (SymTrace's suspicious-node intervention \cite{luan2026symtrace}) \\
 & \textsc{fix-observation} & replace a corrupted/empty $o_k$ (404, OCR failure, swallowed output) by the correct observation \\
 & \textsc{reassign} & let a different agent $a'\neq a_k$ produce step $k$ \\\midrule
$\Delta_{\mathrm{inj}}$ & \textsc{corrupt-output} & rewrite $u_k$ with a context-aware error of mode $m$ (the Aegis/Who\&When~Pro injectors) \\
 & \textsc{corrupt-observation} & perturb $o_k$ (wrong number, truncated page, failing tool) \\
 & \textsc{delete} & drop step $k$ and continue \\
 & \textsc{swap-agent} & route step $k$ to an agent not designed for it \\\bottomrule
\end{tabularx}
\captionof{table}{A starting intervention vocabulary $\Delta$. Each entry is a deterministic transformation; what varies across rollouts is only the continuation.}
\label{tab:vocab}
\end{minipage}\medskip

\subsection{Proposal 1: Counterfactual perturbation prediction on agent trajectories}\label{sec:future:perturb}

\begin{multicols}{2}
\paragraph{Motivation} Every dataset in this survey defines the decisive step \emph{counterfactually}---``the earliest step such that, had it been correct, the run would have succeeded''---but only TraceElephant verifies that counterfactual by re-execution, and the injected datasets (Who\&When~Pro, Aegis, TracerTraj, CORRECT-Error) keep only the injections that produced a failure and discard the ones the system absorbed. Three observations from the raw data make the gap concrete:
\begin{itemize}[leftmargin=1.1em]
\item Same task, same system family, re-run: the blamed agent agrees between Who\&When and TraceElephant in only 18 of 86 paired runs, i.e.\ the root cause of a task failure is run-specific.
\item In LongRCA, 364 of 1,140 reasons state no causal link to the outcome and roughly one in five blames ``redundant/duplicate'' work in runs that still produced a final plan---the label was never tested.
\item Aegis reports that injections flip a successful run into a failure between 34\% (DyLAN) and 77\% (LLM Debate) of the time, so absorption is the common case, yet no dataset records it.
\end{itemize}
Two replay studies confirm the picture from the execution side. SymTrace \cite{luan2026symtrace} re-ran 54 successful MAS executions three times each and saw 39 of them fail at least once, while unguided reruns of failed executions reproduced the failure in only 67.97\% of attempts; in the faulty-agent injections of \citet{huang2024resilience} a reviewing Inspector agent recovered up to 96.4\% of the injected errors. Absorption and run-to-run drift are large effects, not noise. The questions below turn attribution from a human-labelled classification target into an \emph{interventional quantity that can be measured}.

\paragraph{Closest prior work} Four strands already contain parts of the programme.
\begin{itemize}[leftmargin=1.1em]
\item \emph{Replay-based repair.} SymTrace \cite{luan2026symtrace} (\S\ref{ds:symtrace}) records a MAS execution, picks an \emph{anchor} node, serves every LLM and tool call before the anchor from the log and lets everything after it run live---the operational form of $\mathrm{Roll}(\tau_{<k}\oplus s_k')$ of Section~\ref{sec:future:prelim}. Its intervention augments the anchor's request with a category-specific repair instruction ($\hat q_k=q_k\oplus\Delta$, the \textsc{hint} of Table~\ref{tab:vocab}). On SymFail (536 evaluator-confirmed failures of AG2, CrewAI and Magentic-One on WebArena-Verified Hard and AssistantBench) this repairs 20.15\% of failures in one attempt, against 6.90\% for three unguided reruns and 3.7--4.3\% for self-reflection and critic-agent baselines, and selective replay reproduces the original failure in 80.78\% of single attempts versus 67.97\% for a full rerun. SymTrace runs one temperature-0 continuation per intervention, so it observes a single draw from $p_\delta$; it takes the anchor from a symptom detector rather than estimating which step is worth intervening on, and it has no injection side.
\item \emph{Counterfactual attribution.} Causal Agent Replay \cite{shah2026car} states the same do-intervention-then-re-execute estimator with confidence intervals and a Monte-Carlo Shapley extension for interacting steps, but validates it only on synthetic structural causal models. Minimal repair families \cite{li2026mrfr} replay \emph{sets} of events and show that single-event replay misses repairs that are only jointly sufficient (exact family recovery on 90 controlled cases and a 24-case four-agent pilot). EMFA \cite{liao2026emfa} verifies the candidates of an error-propagation model with counterfactual checks and is the current best step-level attributor on Who\&When (+3.45 / +4.40 pp over the previous best on the two subsets); TraceElephant's dynamic setting (\S\ref{ds:traceelephant}) is the same idea with a replayable environment.
\item \emph{Recovery as a measured property.} \citet{huang2024resilience} inject faulty agents into six MAS and find that a hierarchical topology loses 5.5\% of performance versus 10.5--23.7\% for flat ones; the hallucination-cascade study of \citet{jamshidi2026cascade} tracks claim-level errors through three-agent chains and finds net attenuation (normalised hallucination 0.422 $\to$ 0.272); ToolMaze \cite{zhu2026toolmaze} scores a perturbation-recovery rate separately from task success and finds it drops by $\approx$37\% under implicit tool failures. None of these ties recovery to a labelled step of a natural failure.
\item \emph{Infrastructure and training.} AgentRewind \cite{zhuang2026agentrewind} checkpoints agent context and environment state jointly so that an agent can roll back---an implementation of Assumption~\ref{as:replay}; Agent-R \cite{yuan2025agentr} trains agents on trajectories spliced at the first erroneous step with a correct branch, i.e.\ it consumes exactly the (prefix, repaired step) pairs that task T1 would predict.
\end{itemize}
What is missing from all of them is the \emph{quantity}: an effect size per (trajectory, step, intervention), estimated from repeated rollouts, on failed \emph{and} successful natural runs.

With Definitions~\ref{def:effect} and \ref{def:fa} in place, the two questions that motivate the proposal are unambiguous:
\begin{enumerate}[leftmargin=1.4em]
\item \textbf{Repair (Q1).} Given a \emph{failed} $\tau$, a step $k$ and $\delta\in\Delta_{\mathrm{fix}}$: what is $p_\delta(\tau,k)$? Equivalently, which $(k,\delta)$ maximise $\phi$?
\item \textbf{Injection (Q2).} Given a \emph{successful} $\tau$, a step $k$ and $\delta\in\Delta_{\mathrm{inj}}$: what is $1-p_\delta(\tau,k)$, the probability that the perturbation cascades into failure rather than being absorbed?
\end{enumerate}
Neither asks for a binary answer: a single re-run is one Bernoulli draw from $p_\delta$, and Section~\ref{ds:traceelephant} shows that draws from the same prefix land on different failures.
\end{multicols}

\begin{figure}[h]\centering
\begin{tikzpicture}[x=1cm,y=1cm,font=\scriptsize,>=Stealth,scale=0.95,every node/.style={transform shape}]
% prefix
\foreach \i in {1,...,5}{\node[draw=rule,fill=surface,minimum width=0.7cm,minimum height=0.45cm,inner sep=1pt] (s\i) at (\i*0.85,0) {$s_{\i}$};}
\node[draw=critical,fill=critical!8,minimum width=0.7cm,minimum height=0.45cm,inner sep=1pt] (sk) at (6*0.85,0) {$s_k$};
\node[above=1pt of sk,text=critical] {intervene: $\delta(s_k)$};
\draw[decorate,decoration={brace,amplitude=3pt,mirror},rule!80!black] (s1.south west) -- (s5.south east) node[midway,below=3pt,text=ink2]{prefix $\tau_{<k}$, state restored};
% rollouts
\foreach \j/\y/\c/\l in {1/0.9/good/$S=1$, 2/0.3/good/$S=1$, 3/-0.3/critical/$S=0$, 4/-0.9/good/$S=1$}{
  \draw[->,\c!70,line width=0.6pt] (sk.east) .. controls (7.0,0) and (7.4,\y) .. (8.4,\y);
  \foreach \i in {0,1,2}{\node[draw=\c!60,fill=\c!6,minimum width=0.45cm,minimum height=0.32cm,inner sep=0pt] at (8.7+\i*0.55,\y) {};}
  \node[right=2pt,text=\c!70!black] at (9.85,\y) {\l};}
\node[text=ink2,anchor=west] at (10.9,0.3) {$N$ rollouts of $\pi,\mathcal{E}$};
\node[text=ink2,anchor=west] at (10.9,-0.2) {$\hat p_\delta=\tfrac{\#\{S=1\}}{N}$};
\node[text=ink2,anchor=west] at (10.9,-0.7) {$\hat\phi=\hat p_\delta-\hat p_0$};
\end{tikzpicture}
\caption{The measured quantity. The state after the prefix is restored, step $k$ is replaced by $\delta(s_k)$, and $N$ independent continuations give a success rate; the same with the original $s_k$ gives the baseline. The effect is the difference.}
\label{fig:perturb}
\end{figure}

\begin{multicols}{2}
\paragraph{Learning tasks}
\begin{description}[leftmargin=1.4em,style=nextline,itemsep=1pt]
\item[T1 --- Outcome prediction (the surrogate).] Learn $f(\tau_{<k},s_k,\delta)\approx p_\delta(\tau,k)$ \emph{without executing}. Evaluate against Monte-Carlo estimates with Brier score, calibration and AUROC; report separately for $\Delta_{\mathrm{fix}}$ and $\Delta_{\mathrm{inj}}$. Baselines: process-reward-model value estimates \cite{wang2024mathshepherd}, the system's own self-monitoring signals (orchestrator ledger flags precede the labelled step in 24 of 90 Magentic-One traces, \S\ref{ds:traceelephant}), and per-source position priors.
\item[T2 --- Attribution by effect.] Predict $k^\star=\arg\max_k\max_{\delta\in\Delta_{\mathrm{fix}}}\phi(\tau,k,\delta)$ and compare with the human labels of Who\&When, TraceElephant, MP-Bench and LongRCA, and with the single-rollout counterfactual verifiers of EMFA and Causal Agent Replay \cite{liao2026emfa,shah2026car}. Disagreements between effect-based and human labels, adjudicated by rollouts, are the interesting output; MP-Bench's three annotators make it possible to ask whether effect size tracks annotator consensus.
\item[T3 --- Fragility and lock-in.] For successful trajectories, the curve $r(k/T)=\mathbb{E}_{\delta\in\Delta_{\mathrm{inj}}}\,p_\delta(\tau,k)$ measures recoverability as a function of position; the first $k$ at which $r$ drops below a threshold is the \emph{point of no return}, the interventional counterpart of Failure-as-a-Process's $t_{\mathrm{lock}}$ (\S\ref{ds:failure_as_process}). Per-agent and per-mode versions of $r$ quantify which roles and which error modes a system can recover from and whether verifier agents change anything; the topology comparison of \citet{huang2024resilience} (hierarchical versus flat) and SymTrace's rerun regressions \cite{luan2026symtrace} are single points on such curves.
\end{description}

\paragraph{The analogy to perturbation biology} The programme is a trajectory-level version of pooled genetic perturbation screens and their in-silico predictors. Perturb-seq \cite{dixit2016perturbseq} applies a library of CRISPR perturbations to a population of cells and reads out the transcriptional phenotype of each cell; GEARS \cite{roohani2024gears} and scGen \cite{lotfollahi2019scgen} then learn to predict the phenotype of \emph{unseen} perturbations from the measured ones, so that most of the screen can be run in silico. Table~\ref{tab:analogy} maps the objects. The correspondence is useful for borrowing methodology (replicate-based effect estimation, essentiality and buffering as graded rather than binary properties, interaction effects between two perturbations) and for seeing where it breaks: a gene library is finite and fixed, whereas edits to a step are open-ended---hence Assumption~\ref{as:vocab}; cells are independent replicates, whereas rollouts share a prefix and a policy, so $N$ must be chosen per prefix; and a rollout costs tokens rather than reagents, which makes the ``screen'' cheap enough to run exhaustively on small benchmarks.
\end{multicols}

\begin{table}[h]\centering\small
\begin{tabularx}{\linewidth}{@{}>{\RaggedRight}X >{\RaggedRight}X >{\RaggedRight}X@{}}\toprule
Perturbation biology & Agent trajectories & Remark \\\midrule
cell with a genetic background & trajectory prefix $\tau_{<k}$ with restored state & the ``background'' includes the task, the system and the prior errors \\
CRISPR knockout / activation & \textsc{delete}, \textsc{corrupt-}$\ast$ / \textsc{hint}, \textsc{oracle-action} & loss- and gain-of-function interventions on one step \\
phenotype readout (expression, viability) & success predicate $S$, optionally the failure mode of the continuation & binary readout today; a richer readout is the new trajectory itself \\
replicate cells & $N$ independent rollouts & estimates $p_\delta$; biology has thousands of replicates, we have $N\approx 5$--$10$ \\
essential gene & $\alpha$-decisive step & graded by effect size, not by a single annotator \\
genetic buffering / robustness & absorbed injection ($\phi\approx0$ on a successful run) & the data Aegis and Who\&When~Pro discard \\
epistasis (double perturbations) & two-step interventions $(k_1,\delta_1),(k_2,\delta_2)$ & tests MP-Bench's claim of several independent causes; cf.\ minimal repair families \cite{li2026mrfr} \\
in-silico perturbation predictor (GEARS, scGen) & surrogate $f$ (task T1) & replaces most rollouts once calibrated \\\bottomrule
\end{tabularx}
\caption{Correspondence between pooled perturbation screens and interventional failure attribution.}
\label{tab:analogy}
\end{table}

\begin{multicols}{2}
\paragraph{Feasibility with the surveyed data}
\begin{itemize}[leftmargin=1.1em]
\item \emph{Environments.} Assumption~\ref{as:replay} holds for SWE-bench / SWE-Bench~Pro (TraceElephant's SWE-Agent runs, TrajErrBench), ALFWorld and WebShop (AgentErrorBench), the TravelPlanner and VitaBench sandboxes (LongRCA) and $\tau^2$-Bench with its simulated user (TrajErrBench, AgentRx); it fails for live-web GAIA, which should be used only to test the surrogate's transfer. SymTrace's recorded-prefix replay \cite{luan2026symtrace} admits WebArena-Verified and AssistantBench runs in addition, at the price of an unreproducible continuation; AgentRewind's joint context/environment checkpoints \cite{zhuang2026agentrewind} are the engineering needed to restore the state after $\tau_{<k}$ cheaply.
\item \emph{Injectors.} Injectors for $\Delta_{\mathrm{inj}}$ exist (Aegis, Who\&When~Pro, TracerTraj) and only need to be run with the negatives retained.
\item \emph{Budget.} A first dataset of 200 failed and 200 successful trajectories, interventions at the labelled step plus four control positions, two interventions each and $N=8$ rollouts is $\approx$25\,k continuations; at the tail lengths observed here (medians of 10--145 steps, 30\,k--260\,k tokens per full trace) this is of the order of $10^9$ tokens, affordable with small open models and well below the cost of the human annotation it replaces.
\end{itemize}

\paragraph{Difficulties to design for}
\begin{itemize}[leftmargin=1.1em]
\item \emph{Stochasticity}: $p_0$ for a failed run is not zero---a plain re-run can succeed (SymTrace: 6.90\% of failed runs pass within three unguided reruns, and 39 of 54 successful runs fail at least once in three reruns \cite{luan2026symtrace})---so effects must be reported with confidence intervals and the evaluation metric must be probabilistic; a single temperature-0 replay, as in SymTrace, EMFA and the minimal-repair-family search, is one Bernoulli draw.
\item \emph{Exposed downstream errors}: a repair at $k$ can succeed up to step $k+m$ and then fail for a new reason; each rollout should therefore be re-attributed (cheaply, by the surrogate) so that ``no flip'' is separable into ``wrong step'' and ``several independent errors''.
\item \emph{Oracle fixes leak the answer}: \textsc{oracle-action} must be constructed from the ground truth but must not inject it verbatim; \textsc{hint} is the weaker, safer variant and both should be reported.
\item \emph{Vocabulary coverage}: some labelled errors (a wrong team composition before the log begins, 20 of 126 Who\&When Algorithm-Generated labels sit on step~0 for this reason) have no single-step intervention and must be declared out of scope rather than silently scored as zero effect.
\item \emph{Cost concentration}: long-horizon sources (LongRCA) dominate the token budget; stratify the rollout budget by trajectory length rather than by trajectory.
\end{itemize}

\paragraph{What would count as a result}
\begin{itemize}[leftmargin=1.1em]
\item A calibrated surrogate (T1) that predicts, from the prefix alone, which steps of a failed run are worth repairing and which steps of a successful run are load-bearing.
\item Effect-based labels (T2) that agree with human consensus where humans agree and resolve the cases where they do not.
\item Fragility curves (T3) that characterise, for the first time with an interventional definition, how far into a run a multi-agent system can still recover from an error.
\end{itemize}
The first two reuse the benchmarks surveyed here as test sets; the third is new measurement that none of them currently supports.
\end{multicols}

\subsection{Proposal 2: Label-free attribution from the flow of trajectories}\label{sec:future:unsup}

\begin{multicols}{2}
\paragraph{Motivation} Every human-labelled dataset in this survey is a \emph{test} set; the only training-scale resources are injected (Section~\ref{sec:g2}) and carry the positional skew of Fig.~\ref{fig:errpos}. \citet{yeh2026oat} show that step labels can be dispensed with altogether: \textsc{Oat} fits a one-class model of \emph{successful} trajectories---about 100 of them, embedded step by step with an LLM and modelled by a neural controlled differential equation \cite{kidger2020ncde}---and flags the steps of a failed run that deviate from the learned dynamics. On the authors' own tool-use set it reaches F1 0.44 against 0.21 for GPT-4o prompting, on Who\&When 0.23 against 0.15 for GPT-5, and it runs 200--5\,000$\times$ faster. Two features of the result matter more than the model.
\begin{itemize}[leftmargin=1.1em]
\item The signal comes from \emph{selection}, not from the dynamics: the paper's own ablations put an RNN next-step predictor on the same step embeddings close behind the CDE (its Fig.~5, no numbers in the text), and the whole out-of-distribution gain (+0.172 AUROC, at $-$0.028 in-domain) comes from the gated control path that suppresses atypical inputs, not from the differential-equation model; a model of \emph{all} trajectories would simply recover the agent's own policy, under which its mistakes are ordinary, high-probability actions. \note{Why the agent's own policy would have high-probability actions of making mistakes?}
\item Selecting successes already requires an outcome label per trajectory, after which \textsc{Oat} discards the failed half of the data.
\end{itemize}
The direction below asks what the principled estimator is at each level of supervision, down to none, and which errors remain identifiable there. It rests on an observation that holds across the survey: at step level, errors are sparse in \emph{both} outcome classes (Table~\ref{tab:contam}). Successful runs are not clean either---Failure-as-a-Process PASS runs contain 2.3 error events per 22.8 steps, all recovered---so ``success'' is a property of the trajectory, not of its steps.

\paragraph{Three levels of supervision} With the notation of Section~\ref{sec:future:prelim} (trajectory $\tau=(s_1,\dots,s_T)$, success predicate $S$), a corpus can carry
\begin{itemize}[leftmargin=1.1em]
\item \textbf{L2}: the step label $t^\star$ (every dataset of Sections~\ref{sec:g1}--\ref{sec:g4});
\item \textbf{L1$^{+}$}: the outcomes of several \emph{continuations of the same prefix}---rollouts under Assumption~\ref{as:replay}, or the two branches of a fault injection---still without any human label;
\item \textbf{L1}: the outcome $S$ only, from the benchmark verifier, unit tests or a judge---free wherever the task has a ground truth;
\item \textbf{L0}: raw logs, neither $S$ nor $t^\star$.
\end{itemize}
\textsc{Oat} sits between L1 and L0: it uses $S$ to select the training set and then ignores it. The proposal is built at L1 and L1$^{+}$: every trajectory is used, failed or not, no step is labelled, and synthetic trajectories are the natural way to add information.

\begin{assumption}[Sparse errors]\label{as:sparse}
The fraction $\varepsilon$ of erroneous steps is small in every trajectory, failed or not; Table~\ref{tab:contam} puts it between 1\% and 20\% on the surveyed data.
\end{assumption}
\begin{assumption}[Absorption]\label{as:absorb}
A trajectory fails if and only if it contains an error from which the run never recovers; the first such step $t^\star$ determines the outcome, and later steps, right or wrong, are not causal. In the terms of Definition~\ref{def:fa}, $t^\star$ is the lock-in point $k_{\mathrm{lock}}$ (AgentRx's first unrecovered failure, Failure-as-a-Process's $t_{\mathrm{lock}}$), not the decisive step $k^\star$, which precedes it by the fix window---median 1 step, but more than 10 in 239 of 1,184 runs (\S\ref{ds:failure_as_process}). A model built on this assumption is therefore expected to sit late relative to decisive-step labels and is scored against both.
\end{assumption}

\begin{definition}[Step hazard]\label{def:hazard}
Let $h_t=h_\theta(\tau_{<t},s_t)\in[0,1]$ be the probability that step $t$ is the first unrecovered error given its prefix. Under Assumption~\ref{as:absorb} the outcome likelihood is the discrete-time survival model
\begin{align*}
P(S{=}1\mid\tau)&=\textstyle\prod_{t=1}^{T}(1-h_t),\\
P(S{=}0,\,t^\star{=}t\mid\tau)&=h_t\textstyle\prod_{s<t}(1-h_s),
\end{align*}
and the attribution is the posterior
$P(t^\star{=}t\mid S{=}0,\tau)=h_t\prod_{s<t}(1-h_s)\big/\big(1-\prod_{s}(1-h_s)\big)$.
\end{definition}

\emph{Link to Proposal 1.} The survival function $\bar S_t(\tau)=\prod_{s\le t}(1-h_s)$ is the probability that a run which has reached $\tau_{\le t}$ still succeeds, so $h_\theta$ is a learned surrogate of the baseline success rate $p_0(\tau,t{+}1)$ of Definition~\ref{def:effect}, and the online lock-in estimate of U3 below is the point of no return of task T3. The two proposals share one object: the rollouts of Section~\ref{sec:future:perturb} measure it, the model here predicts it from the log.

The levels are different ways of fitting $h_\theta$ to the same likelihood (Fig.~\ref{fig:hazard}).
\begin{itemize}[leftmargin=1.1em]
\item \textbf{L2.} The posterior is supervised directly.
\item \textbf{L1$^{+}$.} A rollout from $\tau_{<k}$ is an ordinary L1 trajectory that shares its first $k-1$ steps with $\tau$ and enters the same likelihood. What it adds is \emph{identification}: with independent trajectories the likelihood only says that some step of a failed run was fatal, whereas two continuations of one prefix with different outcomes place the hazard at or after the branch point, and $N$ continuations give a direct target $\hat p_0$ for $\bar S_{k-1}$---the Monte-Carlo value estimation of process-reward models \cite{wang2024mathshepherd}, applied to agent runs. It is the only step-level supervision that costs no label, and it is exactly the by-product of the screen of Proposal~1.
\item \textbf{L1.} $t^\star$ is marginalised out and $\theta$ maximises $\sum_\tau\log P(S_\tau\mid\tau)$ over \emph{all} trajectories: a failed run contributes ``some step was fatal'', a successful run ``no step was''. This is multiple-instance learning \cite{dietterich1997mil} with the survival link in place of the usual noisy-OR, or equivalently positive--unlabelled learning at step level \cite{elkan2008pu,kiryo2017nnpu}, where the steps of successful runs are (nearly) positive examples of correctness and the class prior is the $\varepsilon$ of Table~\ref{tab:contam}.
\item \textbf{L0.} $S$ is latent as well. The likelihood then needs a model of the steps themselves and becomes a \emph{change-point mixture}: every trajectory follows a shared pre-change dynamics $q_0(s_t\mid\tau_{<t})$---the dominant mode---up to $t^\star\in\{1,\dots,T,\infty\}$, after which a second regime $q_1$ applies; $t^\star$ and $S=\mathbb{1}[t^\star{=}\infty]$ are inferred jointly by EM or by Bayesian change-point detection \cite{adams2007bocpd}.
\end{itemize}
\textsc{Oat} is the special case in which $q_0$ is fitted on known successes and the posterior is replaced by a threshold on reconstruction error.

\paragraph{What synthetic data adds} Besides rollouts, the injectors of Table~\ref{tab:vocab} ($\Delta_{\mathrm{inj}}$, released with Aegis, Who\&When~Pro and TracerTraj) produce from one successful run a failed continuation with a known perturbed step and, as often as not (absorption 23--66\% in Aegis), a successful continuation that contains a step known to be wrong. Both enter the L1 likelihood as ordinary trajectories, so the corpus grows without a label. The pair is also the sharpest test of the thesis: a hazard model must put high posterior on the injected step of the failed branch and low cumulative hazard on the identical step of the absorbed branch, whereas any detector of deviation from the success flow (\textsc{Oat}, the one-class baselines) sees the same step twice and must flag both. Two controls follow. Inject uniformly over positions, because the released injected sets do not (Fig.~\ref{fig:errpos}) and an L1 model would otherwise learn the injector's position prior; and hold out injected modes, so that generalisation across error types is measured rather than memorised. The released injected datasets themselves are used as Tier~C test sets (below), not as training data, for the same reason.
\end{multicols}

\begin{table}[h]\centering\small
\begin{tabular}{@{}llrrr@{}}\toprule
Dataset & Error unit & per trajectory & steps per trajectory & $\varepsilon$ \\\midrule
LongRCA & root step & 1 & 156 & 0.6\% \\
Failure as a Process (FAIL) & $t_{\mathrm{err}}$ / error events & 1 / 2.5 & 41.7 & 2.4\% / 6.0\% \\
Failure as a Process (PASS) & recovered error events & 2.3 & 22.8 & 10.1\% \\
TraceElephant & decisive LLM call & 1 & 27.1 & 3.7\% \\
Who\&When & decisive step & 1 & 22.2 & 4.5\% \\
AgentRx (public subset) & failure events & 4.6 & 44.7 & 10.3\% \\
TRAIL & spans carrying $\ge$1 error & 577 of 4\,626 spans & 31.3 & 12.5\% \\
AgentHallu & responsible step & 1 & 7.5 & 13.3\% \\
MP-Bench ($\ge$2 of 3 votes) & flagged steps & 4.2 & 23.3 & 18.0\% \\\bottomrule
\end{tabular}
\caption{Step-level error rate $\varepsilon$ on the surveyed data (means from the statistics tables of Sections~\ref{sec:g1}--\ref{sec:g4}). Even under the most generous labelling, more than 80\% of the steps of a failed trajectory are correct, and successful trajectories contain recovered errors at a comparable rate.}
\label{tab:contam}
\end{table}

\begin{figure}[h]\centering
\begin{tikzpicture}[x=1cm,y=1cm,font=\scriptsize,>=Stealth]
\def\dx{0.9}
\foreach \i in {1,...,5}{\node[draw=rule,fill=surface,minimum width=0.72cm,minimum height=0.42cm,inner sep=1pt] (s\i) at (\i*\dx,0) {$s_{\i}$};}
\node[draw=critical,fill=critical!10,minimum width=0.72cm,minimum height=0.42cm,inner sep=1pt] (s6) at (6*\dx,0) {$s_{t^\star}$};
\foreach \i in {7,8,9}{\node[draw=critical!50,fill=critical!4,minimum width=0.72cm,minimum height=0.42cm,inner sep=1pt] (s\i) at (\i*\dx,0) {$s_{\i}$};}
\node[right=2pt of s9,text=ink2] {$S=0$};
\draw[decorate,decoration={brace,amplitude=3pt},rule!80!black] (s1.north west) -- (s5.north east) node[midway,above=3pt,text=ink2]{shared flow $q_0$ (dominant mode)};
\draw[decorate,decoration={brace,amplitude=3pt},critical!70!black] (s6.north west) -- (s9.north east) node[midway,above=3pt,text=critical!70!black]{post-change regime $q_1$};
\foreach \i/\h in {1/0.05,2/0.08,3/0.06,4/0.12,5/0.10,6/0.75,7/0.30,8/0.25,9/0.40}{
  \fill[s1!70] (\i*\dx-0.18,-1.35) rectangle (\i*\dx+0.18,-1.35+\h);}
\draw[rule!80!black] (0.5,-1.35) -- (9*\dx+0.45,-1.35);
\node[anchor=east,text=ink2] at (0.4,-1.1) {hazard $h_t$};
\foreach \i/\h in {1/0.053,2/0.080,3/0.055,4/0.104,5/0.076,6/0.520,7/0.052,8/0.030,9/0.036}{
  \fill[critical!70] (\i*\dx-0.18,-2.35) rectangle (\i*\dx+0.18,-2.35+\h);}
\draw[rule!80!black] (0.5,-2.35) -- (9*\dx+0.45,-2.35);
\node[anchor=east,text=ink2] at (0.4,-2.1) {$P(t^\star{=}t\mid S{=}0,\tau)$};
\node[anchor=north west,text=ink2,align=left] at (9*\dx+1.6,0.0) {L2: $t^\star$ observed\\L1: $S$ observed, $t^\star$ marginalised\\L0: $S$ and $t^\star$ latent,\\\phantom{L0: }$q_0$ and $q_1$ fitted jointly};
\end{tikzpicture}
\caption{The hazard model. Each step receives a prefix-conditional hazard $h_t$; the survival product turns the hazards into an outcome probability (what L1 supervises) and, for a failed run, into a posterior over the first fatal step (the attribution). At L0 the same posterior is a change point between a shared dynamics $q_0$ and a post-error regime $q_1$.}
\label{fig:hazard}
\end{figure}

\begin{multicols}{2}
\paragraph{Learning tasks}
\begin{description}[leftmargin=1.4em,style=nextline,itemsep=1pt]
\item[U1 --- Outcome-supervised hazard (L1 and L1$^{+}$).] Fit $h_\theta$ by maximum likelihood on all trajectories of a corpus with outcomes, using the same frozen step embeddings as \textsc{Oat} so that the comparison isolates the estimator. The estimator itself is small: a causal GRU or two-layer transformer over the embeddings with a sigmoid hazard head, a sparsity prior on $\sum_t h_t$ (Assumption~\ref{as:sparse}), and the policy decision of Section~\ref{sec:future:prelim} as the step unit. Steps are discrete events, so no continuous-time machinery is needed---the neural CDE of \textsc{Oat} earns its place for irregularly sampled signals, which agent logs are not; what matters is causality (prefix only) and sparsity. Fit once at L1 and once with the rollout prefixes of L1$^{+}$ added, to measure what synthetic branching buys. Report the attribution posterior against the set-valued labels (MP-Bench consensus ranking with nDCG, AgentRx failure events, TRAIL spans, $t_{\mathrm{err}}$ and $t_{\mathrm{lock}}$ of Failure-as-a-Process) and exact / $\pm k$ step accuracy on Who\&When and TrajErrBench; report the calibration of $P(S{=}0\mid\tau)$, which the model yields for free.
\item[U2 --- Fully unsupervised change-point model (L0).] Same corpora, outcomes hidden. The question is not whether it beats U1---it cannot---but \emph{which error classes are identifiable from the dominant mode alone}. The generator-likelihood argument predicts that errors visible only in the action stream are not (a hallucinated summary is a typical summary), whereas errors that change the \emph{observation} stream are: tool exceptions, empty pages, repeated queries (MAST FM-1.3), stalled ledgers, premature termination. Measure recall per category on the taxonomy datasets (TrajErrBench modules, Who\&When~Pro families, TRAIL and AgentRx categories), fitting $q_0$ on observations and on actions separately.
\item[U3 --- Prefix monitor.] Because $h_t$ depends on the prefix only, the cumulative hazard $1-\prod_{s\le t}(1-h_s)$ is an online estimate of lock-in. Evaluate its lead time against $t_{\mathrm{lock}}$ (Failure-as-a-Process's prompted monitor reaches 82\% precision at 18\% recall with median lead time 0) and against the orchestrator's own ledger flags in Magentic-One traces (\S\ref{ds:traceelephant}).
\end{description}

\paragraph{Evaluation suite} Three tiers, so that the method is scored where its target is measurable, where it must transfer, and where labels are exact.
\begin{itemize}[leftmargin=1.1em]
\item \emph{Tier A, measured} (train and test in the same replayable system; rows 1--5 of Table~\ref{tab:engine}). $h_\theta$ is fitted on self-collected runs and scored on the released human labels of TrajErrBench, AgentErrorBench and the SWE-Agent / SWE-Bench~Pro traces; on a subset the rollout screen of Proposal~1 supplies $\hat p_0$ curves, so $\bar S_t$ is scored directly by Brier score and calibration---the only tier where the object itself, not a proxy, is the reference.
\item \emph{Tier B, transfer} (log-only; row 6). Frozen model, new systems: Who\&When (both subsets), MP-Bench (nDCG@5 against the consensus ranking), LongRCA (long horizon, $\pm5$ window), AgentRx (set-F1 over failure events) and AgentHallu, whose 250 clean runs allow an L1 fit from the dataset's own logs.
\item \emph{Tier C, exact labels} (synthetic). Who\&When~Pro text split, CORRECT-Error, TracerTraj test, and our own injections with the absorbed branches retained; results stratified by injected position.
\end{itemize}
Metrics: exact and $\pm k$ step accuracy; distance to $t_{\mathrm{err}}$ and to $t_{\mathrm{lock}}$ reported separately (Assumption~\ref{as:absorb} predicts a late bias against the former); set metrics where the label is a set; lead time for U3; calibration of $P(S{=}0\mid\tau)$. Baselines, all on the same frozen encoder: position priors (first step, last decision step, per-source histogram---the last decision step alone is inside the tolerance span of 16.9\% of SearchAuditBench labels and the last fifth holds 52\% of CORRECT-Error labels); All-at-Once prompting; \textsc{Oat} re-implemented; Isolation Forest and one-class SVM on step embeddings; and a value regression trained on rollouts only \cite{wang2024mathshepherd}, i.e.\ L1$^{+}$ without L1. Robustness: outcome labels flipped at rates up to the 12.9\% verifier--judge disagreement observed in HORIZON, and judge-only against verifier-only outcomes on the same runs.

\paragraph{Pitfalls specific to this data}
\begin{itemize}[leftmargin=1.1em]
\item \emph{The last-step shortcut.} The final answer reveals the outcome, so an L1 model will put hazard on it; the injected datasets show how strong that prior is (median relative position 0.83 in Who\&When~Pro, 0.95 in CORRECT-Error). Mask the last $k$ steps from the hazard input and report accuracy with the final step excluded.
\item \emph{Task difficulty as a confound.} Hard tasks produce both failures and unusual steps, so an outcome model can learn ``this task is hard'' instead of ``this step is wrong''. A matched design---the same task run to both outcomes---removes it; the 89 tasks $\times$ 21 systems of Section~\ref{ds:failure_as_process} are the template.
\item \emph{Recovered errors.} The survival likelihood tolerates them (a success only requires that no error was fatal); the plain multiple-instance assumption ``success bags are clean'' does not. Use the positive--unlabelled formulation with the prior from Table~\ref{tab:contam} when the any-error target is wanted.
\item \emph{Identifiability and model selection at L0.} A change-point mixture can converge to a split that is not error-related (planning phase versus execution phase, say). Use a small L2 set for model selection and calibration only, never for training, and say so.
\item \emph{Anomaly is not causality.} Both \textsc{Oat} and U2 flag deviation; only the counterfactual rollouts of Section~\ref{sec:future:perturb} can confirm that a flagged step is decisive. The two directions close a loop: the hazard posterior is a cheap proposal distribution for where to intervene, and the rollouts return labels that re-fit $h_\theta$ at L2.
\end{itemize}

\paragraph{Feasibility with the surveyed data}
\begin{itemize}[leftmargin=1.1em]
\item \emph{Released corpora with both outcomes.} AgentHallu (443 hallucinated / 250 clean), HORIZON (2\,498 / 1\,284 by judge), TRAIL (148 runs not selected by outcome) and Failure-as-a-Process (1\,184 / 610, annotations only---the CodeTracer trajectories must be re-joined). The core MAS benchmarks release failures only (TraceElephant withholds its 160 successful runs; Who\&When~Pro and Aegis do not ship their successful seeds).
\item \emph{Collection engine.} Everything else is collected, not downloaded: Table~\ref{tab:engine} lists the systems of the survey whose code and environments are public and replayable, with the verifier that supplies $S$ and the released labels that serve as test sets on them. One thousand runs of one system cost about $1{,}000\times30$ steps $\times\,5$\,k tokens $\approx1.5\times10^{8}$ tokens with an open model, the outcome is free, and the same infrastructure (state restore, LLM-call cache) is what Proposal~1 needs for its rollouts, so the proposals share one data engine.
\item \emph{Evaluation.} The set-labelled datasets for U1/U2 and the three timestamps of Failure-as-a-Process for U3 exist already; no new annotation is needed.
\item \emph{Budget.} One forward pass per step for the embeddings: $\approx$2\,k trajectories $\times$ 30\,k tokens $\approx6\times10^{7}$ tokens, three to four orders of magnitude below the perturbation screen of Section~\ref{sec:future:perturb}; the hazard head trains in minutes.
\end{itemize}

\paragraph{What would count as a result}
\begin{itemize}[leftmargin=1.1em]
\item An L1 estimator that matches or exceeds \textsc{Oat} and prompting on the natural benchmarks with the same encoder and no step labels, with a calibrated outcome probability as a by-product.
\item A measured gain from synthetic branching: the L1$^{+}$ fit beats the L1 fit on Tier~A at equal label cost (zero), and the hazard model separates injected from absorbed branches where one-class detectors cannot.
\item An identifiability map for L0: the error categories that can and cannot be recovered from the dominant mode, which tells practitioners what a label-free monitor will and will not catch.
\item A prefix monitor with positive median lead time to $t_{\mathrm{lock}}$.
\end{itemize}
\end{multicols}

\medskip\noindent\begin{minipage}{\linewidth}\centering\scriptsize
\setlength{\tabcolsep}{3pt}
\begin{tabularx}{\linewidth}{@{}>{\RaggedRight\hsize=1.15\hsize}X >{\RaggedRight\hsize=0.8\hsize}X >{\RaggedRight\hsize=0.65\hsize}X >{\RaggedRight\hsize=1.2\hsize}X >{\RaggedRight\hsize=1.2\hsize}X@{}}\toprule
System family (public agent code) & State restore for $\mathrm{Roll}$ & Verifier $S$ & Released labels usable as test set & Remark \\\midrule
1\ \ $\tau^2$-Bench airline / retail: tool-calling agent + simulated user (TrajErrBench, AgentRx) & reset of domain DB and user simulator; seeded & DB-state constraints & TrajErrBench 400 ($\kappa$ 0.91), AgentRx 29 & cheapest: 29 messages, 20\,k characters per run \\
2\ \ ALFWorld, WebShop: AgentDebug's modular ReAct agent & deterministic text simulators & task reward & AgentErrorBench 100 + 50 & 30-step cap; observations repeat the page \\
3\ \ SWE-bench Verified / Pro, Terminal-Bench 2: SWE-Agent with TraceElephant's tracing middleware, OpenHands, mini-SWE-agent, Terminus~2 & container snapshot (Harbor) & hidden tests & TraceElephant 44, TrajErrBench 86, Failure-as-a-Process $t_{\mathrm{err}}/t_{\mathrm{lock}}$ for 1,184 runs once re-joined with the CodeTracer release & long: 120--300 steps, 200\,k characters; dominates the budget \\
4\ \ MASLab systems (LLM Debate, MacNet, AgentVerse, DyLAN, SmolAgents, Magentic-One) on GSM8K / MATH / HumanEval / MMLU-Pro: Aegis wrappers and injection factory; MetaGPT / AFlow runs of TracerTraj & no environment state; LLM-call cache, temperature 0 & answer match, unit tests & Aegis 600 (agent + mode, no step), TracerTraj test 790 (synthetic) & injector released; the natural setting for $\Delta_{\mathrm{inj}}$ with absorbed branches kept \\
5\ \ WebArena (Verified), AgentBench-OS, BIRD: HORIZON agents; LongRCA's planner--surfer--verifier controller & docker reset per task & environment verifier & HORIZON 294 judge step labels, LongRCA WebArena 177 & HORIZON agents public; LongRCA system code not released \\
6\ \ Live-web GAIA / AssistantBench / BrowseComp: Magentic-One, CaptainAgent, smolagents, deep-search scaffolds & none (recorded-prefix replay only) & answer match & Who\&When 184, TraceElephant 176, MP-Bench 289, SymFail 536, TRAIL 117, AgentHallu 36, SearchAuditBench 1,243 & Assumption~\ref{as:replay} fails: transfer tests only \\\bottomrule
\end{tabularx}
\captionof{table}{The data engine shared by the three proposals: systems of the survey on which unlabelled trajectories can be collected with replay, the verifier that gives the outcome for free, and the released labels that score the result on them.}
\label{tab:engine}
\end{minipage}\medskip

\subsection{Proposal 3: Attribution-guided policy improvement}\label{sec:future:train}

\begin{multicols}{2}
\paragraph{Motivation} The datasets of this survey are built to score \emph{attributors}. The obvious consumer of a good attributor is the agent that produced the failure: a label $(a^\star,k^\star)$ says which parameters to move and from which state to try again. A literature search on this direction (notes in \code{notes/attribution\_as\_training\_signal.md}) gives three findings.
\begin{itemize}[leftmargin=1.1em]
\item \emph{The mechanism exists, under other names.} The strongest 2026 agent-RL methods locate a ``pivotal'', ``critical'' or ``irrecoverable'' step of a failed rollout and use it to restart rollouts there \cite{tian2026mirl,guo2026pivoarl,shi2026r3l,liu2026sprout,liang2026elpo}, to reweight the advantage at that step \cite{zhu2026fault,huang2026etro,jung2026prover,agentevolver2025,zhu2026lvrl} or to build preference pairs there \cite{li2026cso,yuan2025agentr,steca2025,cuarecovery2026}. None of them cites a failure-attribution benchmark as a source of signal, and none uses an attributor trained on one.
\item \emph{The benchmarks never close the loop.} Every downstream experiment inside an attribution paper is inference-time: AgenTracer feeds its tracer's reasoning back as a prompt for three rounds (gains of 4.8--14.2\%) \cite{agentracer2026}, AgentDebug re-rolls from the critical step with guidance \cite{agentdebug2026}, SearchAuditor resumes at the critical step with a repair directive \cite{searchauditor2026}, SymTrace replays one node \cite{luan2026symtrace}. No weights change. Aegis, AgentRx, TraceElephant, Who\&When~Pro and MP-Bench have no downstream experiment, and the only multi-agent blueprint, first-error localisation followed by preference pairs on the blamed agent, is a position paper without experiments \cite{yang2025blame}.
\item \emph{Attribution accuracy is never related to training value.} Nobody has asked whether a step metric from Sections~\ref{sec:g1}--\ref{sec:g4} predicts how much the agent improves when the label is used. The one direct measurement is negative: on 150 SWE-bench Verified runs, LLM-judge ``first mistake'' labels fall where the success probability does \emph{not} drop in 66\% of cases, only 7 of 188 judge answers name the decisive step, and restarting an agent from the judge's step changes expected success by $-0.03$ to $-0.08$, whereas restarting from the certified point of no return raises it from 0.44 to 0.53, and from the measured restart point to 0.57 \cite{shen2026epv}. SearchAuditor's own ablation adds that, for Kimi-K2.6, a generic hint at the same predicted step repairs fewer failures (5.6\%) than an unguided retry (9.4\%), against 17.4\% with the auditor's directive \cite{searchauditor2026}: the location is necessary, not sufficient.
\end{itemize}
The pattern shared by the methods that work is \emph{propose, verify, then credit}: a judge or the policy proposes a step, a few current-policy continuations confirm that the outcome actually turns there, and only then does the step shape credit or data \cite{liang2026elpo,jung2026prover,li2026cso,chu2026agentgrad}. That verification step is the rollout screen of Proposal~1 with $N$ between 1 and 8. The proposal below makes the link explicit and asks the question the literature skips.

\paragraph{Setting} The agent is a trainable policy $\pi_\theta$ (single- or multi-agent, each agent $a\in\mathcal{A}$ owning a parameter subset $\theta_a$) run on tasks with a success predicate $S$; the notation is that of Section~\ref{sec:future:prelim}. A group-relative outcome-only update (GRPO and its agentic variants \cite{feng2025gigpo}) applies to every step of a failed trajectory the same negative advantage $A_\tau<0$:
\[g^{\mathrm{unif}}_\tau=A_\tau\sum_{t=1}^{T}\nabla_\theta\log\pi_\theta(u_t\mid\tau_{<t}).\]
Under Assumption~\ref{as:absorb} only the steps in the error set are causal, so this estimator pushes down $T-|E_\alpha(\tau)|$ correct decisions per failed run; with the error rates of Table~\ref{tab:contam} ($\varepsilon$ between 1\% and 20\%) that is 80--99\% of its negative gradient mass. An \emph{attributor} $\mathcal{A}$ maps a trajectory to a non-negative weight vector $\hat w_\tau\in\mathbb{R}^T$, an estimate of where the causal mass lies: a point estimate $\hat w_\tau=\mathbb{1}[t=\hat k]$, a posterior $\hat w_\tau(t)=P(t^\star{=}t\mid S{=}0,\tau)$ from the hazard model of Definition~\ref{def:hazard}, or a graded profile $\hat w_\tau\propto\hat e_\tau$ from the screen of Definition~\ref{def:fa}.

\begin{definition}[Attribution-weighted update]\label{def:attrgrad}
Given $\hat w_\tau$ and a mixing weight $\lambda\in[0,1]$, the attribution-weighted policy gradient of a failed trajectory is
\[g^{\hat w}_\tau=A_\tau\sum_{t=1}^{T}\big[(1-\lambda)+\lambda\,T\hat w_\tau(t)/\|\hat w_\tau\|_1\big]\nabla_\theta\log\pi_\theta(u_t\mid\tau_{<t}),\]
so that $\lambda=0$ recovers $g^{\mathrm{unif}}_\tau$ and $\lambda=1$ with the true $w_\tau=\mathbb{1}[t=k^\star]$ concentrates the whole penalty on the decisive step. Successful trajectories are unchanged.
\end{definition}

The bias of $g^{\hat w}$ relative to the ideal $g^{w}$ is governed by the attributor's precision at the chosen tolerance, which is exactly what the step metrics of this survey measure; its variance reduction relative to $g^{\mathrm{unif}}$ is governed by $1-\varepsilon$. The trade-off is therefore a function of a quantity the benchmarks already report, and the proposal is to measure the function.

\paragraph{Three entry points} The weight vector can enter training in three places, each with a different infrastructure requirement (Fig.~\ref{fig:entry}); existing methods occupy them as listed in Table~\ref{tab:entry}.
\begin{description}[leftmargin=1.4em,style=nextline,itemsep=1pt]
\item[R --- Restart.] Rollouts for the next update start from $\tau_{<k}$ with $k\sim\rho_\tau\propto\hat w_\tau$ instead of from the task prompt. This reshapes the state distribution toward decisions that matter and lets a group of continuations share a prefix, so GRPO's within-group comparison becomes a step-level comparison at $k$. It needs Assumption~\ref{as:replay}: MiRL keeps only the 34\% of prefixes whose commands replay deterministically \cite{tian2026mirl}, SPROUT snapshots the file system and context after every turn and keeps all of them \cite{liu2026sprout}.
\item[C --- Credit.] The update of Definition~\ref{def:attrgrad}, optionally with \emph{truncation}: no dense supervision after $\hat k$, because later steps are conditioned on a state the current policy would not reach \cite{huang2026etro,guo2026pivoarl}. It needs only the log and an attributor.
\item[D --- Data.] A preference pair or a spliced correction at $\hat k$: the policy's action against an alternative from which a continuation succeeds \cite{li2026cso}, or the failed prefix followed by a successful sibling branch \cite{yuan2025agentr}. It needs an alternative action and a way to check it, i.e.\ one rollout from $\tau_{<\hat k}\oplus\delta(s_{\hat k})$ with $\delta\in\Delta_{\mathrm{fix}}$ of Table~\ref{tab:vocab}.
\end{description}
In a multi-agent system $\hat w$ is indexed by (agent, step) and a fourth entry point opens: \textbf{M --- culprit-only update}, the gradient is applied to $\theta_{a^\star}$ and the advantages of the other agents are masked, the weight-space analogue of what intervention-guided prompt optimisation already does \cite{chu2026agentgrad}. The RL literature on LLM teams instead uses dense per-agent credit, from a coach that scores every action \cite{li2026mappa}, from leave-one-out masking \cite{li2026sharp} or from exact counterfactual continuation of every message \cite{chen2026c3}, and only on two- to three-agent mathematics; nothing uses (agent, step) blame on a GAIA-style system.
\end{multicols}

\begin{figure}[h]\centering
\begin{tikzpicture}[x=1cm,y=1cm,font=\scriptsize,>=Stealth]
\def\dx{0.85}
\foreach \i in {1,...,4}{\node[draw=rule,fill=surface,minimum width=0.68cm,minimum height=0.4cm,inner sep=1pt] (s\i) at (\i*\dx,0) {$s_{\i}$};}
\node[draw=critical,fill=critical!10,minimum width=0.68cm,minimum height=0.4cm,inner sep=1pt] (s5) at (5*\dx,0) {$s_{\hat k}$};
\foreach \i in {6,7}{\node[draw=rule,fill=surface,minimum width=0.68cm,minimum height=0.4cm,inner sep=1pt] (s\i) at (\i*\dx,0) {$s_{\i}$};}
\node[right=2pt of s7,text=ink2] {failed run, $S=0$};
\foreach \i/\h in {1/0.04,2/0.06,3/0.05,4/0.12,5/0.58,6/0.07,7/0.04}{\fill[s1!70] (\i*\dx-0.17,-1.0) rectangle (\i*\dx+0.17,-1.0+\h);}
\draw[rule!80!black] (0.5,-1.0) -- (7*\dx+0.4,-1.0);
\node[anchor=east,text=ink2] at (0.45,-0.75) {attribution $\hat w_\tau(t)$};
\coordinate (p) at (5*\dx,-1.08);
\node[draw=s3,fill=s3!8,rounded corners=2pt,inner sep=3pt,text width=4.6cm,align=left] (R) at (5*\dx-5.3,-2.55) {\textbf{R --- restart.} Next rollouts start from $\tau_{<k}$ with $k\sim\rho_\tau\propto\hat w_\tau$. Needs state restore.};
\node[draw=s2,fill=s2!8,rounded corners=2pt,inner sep=3pt,text width=4.6cm,align=left] (C) at (5*\dx,-2.55) {\textbf{C --- credit.} Advantage of step $t$ scaled by \mbox{$(1-\lambda)+\lambda T\hat w_\tau(t)$}, truncated after $\hat k$. Needs only the log.};
\node[draw=s7,fill=s7!8,rounded corners=2pt,inner sep=3pt,text width=4.6cm,align=left] (D) at (5*\dx+5.3,-2.55) {\textbf{D --- data.} Pair $(u_{\hat k},\delta(u_{\hat k}))$, kept if a rollout from $\delta(u_{\hat k})$ succeeds. Needs one verifying rollout.};
\draw[->,ink3] (p) -- (R.north);
\draw[->,ink3] (p) -- (C.north);
\draw[->,ink3] (p) -- (D.north);
\end{tikzpicture}
\caption{Entry points of an attribution estimate $\hat w_\tau$ into agent training. All three consume the same object and differ in what they need from the environment. In a multi-agent system the same vector, indexed by agent and step, also selects whose parameters are updated (entry point M).}
\label{fig:entry}
\end{figure}

\begin{table}[h]\centering\scriptsize
\setlength{\tabcolsep}{3pt}
\begin{tabularx}{\linewidth}{@{}l >{\RaggedRight\hsize=1.05\hsize}X l >{\RaggedRight\hsize=0.95\hsize}X l@{}}\toprule
Method & Attributor (who proposes $\hat k$) & Verified? & Entry & Gain vs.\ outcome-only \\\midrule
MiRL \cite{tian2026mirl} & 27B judge, clustered into 8 misbehaviours & replay only & R + judge-gated reward & TB-Lite 51.8 vs 43.4 \\
PivoARL \cite{guo2026pivoarl} & the policy's own reflection & no & R + C (prefix credit) & 83.9 vs 73.9 (GiGPO) \\
R$^3$L \cite{shi2026r3l} & the policy's own diagnosis & no & R + C (prefix mask) & ALFWorld 0.72 $\to$ 0.93 \\
SPROUT \cite{liu2026sprout} & policy as reviewer, with test diagnostics & no & R (snapshots) & TB-2.1 64.0 vs 39.3 \\
ELPO \cite{liang2026elpo} & binary-search rollouts: first irrecoverable step & MC & R + C (suffix clip) & 60.7 vs 51.2 \\
FAULT \cite{zhu2026fault} & self-diagnoser with step and quote, evidence-checked & quote match & C (priced penalties) & ALFWorld 91.0 vs 49.3 \\
ETRO \cite{huang2026etro} & 27B writer: earliest unrecovered decisive error & no & C (truncated distillation) & WebShop 85.8 vs 73.7 \\
ProVer \cite{jung2026prover} & agentic judge proposes a segment & MC, $N{=}8$ & C (positive only) & $+7$ to $+10\%$ rel. \\
AgentEvolver \cite{agentevolver2025} & judge labels every step good/bad & no & C ($\alpha{=}0.1$--$0.2$) & 47.6 vs 41.3 \\
LVRL \cite{zhu2026lvrl} & frozen LLM value function, credit = value drop & no & C & up to $+13.6$ pp \\
CSO \cite{li2026cso} & PRM candidates, expert alternative & MC, policy continuation & D (DPO at $\hat k$) & GAIA 49.5 vs 38.9 \\
Agent-R \cite{yuan2025agentr} & the policy's own first-error judgement & no & D (splice) & 70.7 vs 62.4 (splice at end) \\
CUA-Recovery \cite{cuarecovery2026} & gold per-step lineage, human-checked & yes & D (recovery SFT) & 35.2 vs 23.6 (token-matched) \\
MAPPA \cite{li2026mappa} & LLM coach scores every (agent, action) & no & M (dense) & no shared-reward ablation \\
C3 \cite{chen2026c3} & counterfactual continuation of every message & MC & M (dense) & MATH500 82.8 vs 74.5 \\
AgentGrad \cite{chu2026agentgrad} & sequential intervention: first agent whose fix flips $S$ & yes & M (prompt only) & $+11.8$ vs $+9.2$ (GEPA) \\\bottomrule
\end{tabularx}
\caption{Existing methods placed on the entry points of Fig.~\ref{fig:entry}. ``Verified'' states whether the proposed step is checked against the outcome before it is used. Numbers are the papers' own and are not compute-matched across rows; see \code{notes/lit\_training\_signal/} for the per-paper caveats.}
\label{tab:entry}
\end{table}

\begin{multicols}{2}
\paragraph{The attributor ladder} The experiments below hold the training recipe fixed and vary only $\mathcal{A}$, ordered by cost and by how much of the outcome it has seen:
\begin{enumerate}[leftmargin=1.4em,itemsep=1pt]
\item \emph{position priors}: last decision step, per-source histogram (Section~\ref{sec:future:unsup});
\item \emph{outcome variance, no attribution}: start states where the reference policy sometimes fails \cite{pivotrl2026};
\item \emph{the policy's own reflection} \cite{guo2026pivoarl};
\item \emph{prompted judge}: the all-at-once protocol of Who\&When with a frontier model;
\item \emph{trained attributors}: an 8B model trained on TracerTraj, Aegis or Who\&When~Pro with the AgenTracer recipe \cite{agentracer2026}, and the hazard model of Proposal~2 fitted at L1 on the same system's runs;
\item \emph{rollout-verified}: any of 3--5 followed by the screen of Proposal~1 with $N\le8$ at the proposed step and its neighbours \cite{liang2026elpo,jung2026prover};
\item \emph{oracle}: the rollout-measured $k^\star$ with large $N$, affordable only on a test subset.
\end{enumerate}
Each rung is scored twice: by the attribution metrics of this survey (exact and $\pm k$ step accuracy, nDCG against MP-Bench votes) on held-out labels of the same system, and by its training value.

\begin{definition}[Training value of an attributor]\label{def:trainval}
For a fixed recipe, budget $B$ of environment steps and seeds, let $\pi^{\mathcal{A}}_B$ be the policy after training with weights from $\mathcal{A}$ and $\pi^{\mathrm{unif}}_B$ the outcome-only policy at the same budget. The training value is $V_B(\mathcal{A})=J(\pi^{\mathcal{A}}_B)-J(\pi^{\mathrm{unif}}_B)$, with $J$ the held-out success rate; the \emph{accuracy-to-value curve} is $V_B$ plotted against the attributor's step accuracy at the tolerance used by the recipe.
\end{definition}

\paragraph{Experiments}
\begin{description}[leftmargin=1.4em,style=nextline,itemsep=1pt]
\item[A1 --- The accuracy-to-value curve.] Recipe: entry points R and C (restart from $\hat k$, prefix-only credit, truncation after $\hat k$) on a GiGPO-style learner \cite{feng2025gigpo}, 4B policy, matched rollout budget. Systems: rows 1--3 of Table~\ref{tab:engine} (ALFWorld and WebShop with the AgentDebug agent, $\tau^2$-Bench, SWE-bench with a snapshotting harness), each with released human step labels (AgentErrorBench, TrajErrBench, TraceElephant) as the accuracy reference. Every rung of the ladder is trained, plus \emph{synthetically degraded} attributors: the oracle label shifted forward by $d$ steps toward the symptom and the oracle mixed with random steps at precision $p$. The output is the curve of Definition~\ref{def:trainval}, its slope, and the precision below which attribution \emph{hurts} relative to outcome-only training. The evidence of \cite{shen2026epv} predicts that rung 4 sits below rung 2 on SWE-bench and that the shift $d$, not the precision $p$, is the damaging error.
\item[A2 --- Verified negative credit.] Entry point C alone, so that nothing is restarted. Compare (i) unverified judge weights \cite{agentevolver2025}, (ii) evidence-checked self-blame \cite{zhu2026fault}, (iii) the hazard posterior of Proposal~2, and (iv) the rollout-verified profile $\hat e_\tau$ of Proposal~1 used as $\hat w$, each at fixed $\lambda$ and at $\lambda$ annealed toward 0. The quantity watched besides $J$ is \emph{reward hacking}: the attributor's own score rising while $J$ falls, which is documented for fixed-weight process rewards \cite{liu2026vppo,choudhury2025agentprm,zhao2026credithurts}. ProVer uses only positive segments \cite{jung2026prover}; this experiment is its negative mirror on natural failures.
\item[A3 --- Culprit-only updates in a multi-agent system.] Entry point M on the MASLab systems of row 4 (Aegis wrappers, verifiable answers) and, for transfer, Magentic-One on GAIA (row 6, log-only). Blame from the trained attributor (rung 5) and from counterfactual message replay \cite{chen2026c3}. Conditions: shared reward to all agents, per-agent leave-one-out credit \cite{li2026sharp}, repair-every-agent data \cite{zhao2025sirius}, culprit-only gradient, and culprit-only preference pairs \cite{yang2025blame}. Scored by $J$ and by the agent-level accuracy of the blame on Who\&When-style held-out labels of the same system (Aegis-Bench for the MASLab systems).
\item[A4 --- Closing the loop.] After each A1--A3 run, the failures of the trained policy are re-attributed and re-labelled with the failure-mode taxonomies of Section~\ref{sec:g3} (MAST modes through the Aegis labeller, AgentErrorTaxonomy). Two numbers: how much the targeted modes shrink, which MiRL reports as 43--57\% on unseen prefixes \cite{tian2026mirl}, and how much the other modes grow. The second is the ``whack-a-mole'' risk no paper reports; the taxonomy datasets of this survey are the measuring instrument.
\end{description}

\paragraph{Pitfalls to design for}
\begin{itemize}[leftmargin=1.1em]
\item \emph{Symptom shift.} Judge labels drift to the step where the failure becomes visible \cite{shen2026epv}; injected datasets place labels late (Fig.~\ref{fig:errpos}); a policy's own reflection is bounded by its own competence. Restart entry points are the most exposed, since a restart after the point of no return explores a dead state. Verify ($N\le8$) or anneal $\lambda$; never use rung 3 or 4 alone for R.
\item \emph{Reward hacking.} Any score the policy can raise without succeeding will be raised: VPPO's policy split steps to collect prefix reward \cite{liu2026vppo}, AgentPRM's process reward rose while success fell from 82\% to 70\% \cite{choudhury2025agentprm}. Guards used in the literature: conjunction with the outcome ($r=r_o\wedge r_e$ \cite{tian2026mirl}), a small or decaying weight \cite{agentevolver2025,zhao2026credithurts}, a penalty budget capped below the success--failure gap and no gradient through the diagnosis tokens \cite{zhu2026fault}.
\item \emph{Baselines that get the gain without attribution.} Variance-selected start states \cite{pivotrl2026}, recovery training forked from the \emph{end} of failed runs with no localisation (Terminal-Bench 34.8 vs 28.0 \cite{reterm2026}), implicit step credit from repeated states \cite{feng2025gigpo}, and SymTrace's last-node and random-node protocol \cite{luan2026symtrace}. A1 includes all of them as rungs or controls.
\item \emph{Off-policy prefixes.} A corrected or teacher-written step inside a prefix is off-policy for the learner; RTI gates it by perplexity \cite{chen2026rti}, R$^3$L strips the guidance before training \cite{shi2026r3l}. Report the fraction of training tokens that are off-policy.
\item \emph{State restore.} Deterministic replay loses two thirds of terminal prefixes \cite{tian2026mirl}; per-turn snapshots lose none \cite{liu2026sprout}. Row 3 of Table~\ref{tab:engine} should use snapshots, and live-web rows cannot use entry point R at all.
\item \emph{Compute matching.} Several headline gains in Table~\ref{tab:entry} come from baselines trained from the raw backbone or with fewer rollouts; A1--A3 match environment steps, not wall-clock, and report both.
\end{itemize}

\paragraph{Feasibility with the surveyed data}
\begin{itemize}[leftmargin=1.1em]
\item \emph{Environments and labels.} Rows 1--4 of Table~\ref{tab:engine} supply replayable systems with released step labels on the same system (AgentErrorBench 200, TrajErrBench 486, TraceElephant's 44 SWE-Agent traces) and Aegis-Bench's (agent, mode) labels for the multi-agent case. MiRL releases code, checkpoint and trajectories under Apache~2.0 \cite{tian2026mirl}; PivoARL and C3 release code.
\item \emph{Attributors.} AgenTracer-8B's weights are not released (\S\ref{ds:tracertraj}), so rung 5 is retrained from TracerTraj (3,998), Aegis (9,533) and Who\&When~Pro text (6,257), all ungated; the hazard model of Proposal~2 is fitted on the runs A1 collects anyway.
\item \emph{Budget.} One RL run at MiRL's scale is 160 updates $\times$ 8 prompts $\times$ 16 rollouts $\approx 2\times10^4$ rollouts; A1 with seven rungs, four degraded variants, three systems and three seeds is $\approx100$ runs, $2\times10^6$ rollouts, between $10^{10}$ and $10^{11}$ tokens with a 4B policy depending on the system (MASLab and $\tau^2$ runs are 5--20\,k tokens, SWE-bench runs up to 50\,k). The verification rollouts of rungs 6--7 add $N\le8$ continuations per failed trajectory, a small fraction of that. A3 on MASLab systems is cheaper (median 10 turns, 5\,k tokens).
\item \emph{Shared infrastructure.} State restore and the LLM-call cache are the same components as Proposal~1's screen; the verified $(\tau,k,\delta,\hat p_\delta)$ tuples that A2 produces are training data for Proposal~2 at level L1$^{+}$, so the proposals share one engine.
\end{itemize}

\paragraph{What would count as a result}
\begin{itemize}[leftmargin=1.1em]
\item A measured accuracy-to-value curve, with the attribution precision below which attribution-weighted training is worse than outcome-only training, and the finding of whether symptom shift or random error is the damaging failure of an attributor.
\item Verified negative credit at the decisive step that beats unverified judge credit at equal budget and does not exhibit the score-up, success-down signature.
\item A culprit-only update that beats the shared reward on a multi-agent system with blame from an attributor trained on the datasets of this survey, which would be the first use of those datasets to improve an agent.
\item A closed-loop measurement showing the targeted failure modes shrink without the others growing, which is the practical meaning of ``learning from failure attribution''.
\end{itemize}
\end{multicols}

